\section{From Program to Process}
\label{sec:program}

\subsection{A Program Is Not a Process}

``What is a program?'' is less obvious than it looks. A program is a piece of code---but \emph{which} code? Any of the following layers can be called ``the program'':
\begin{itemize}
  \item high-level languages: C, C++, Java, \dots;
  \item intermediate languages: Java bytecode, LLVM IR, .NET CIL, \dots;
  \item low-level languages: assembly;
  \item machine code.
\end{itemize}
A process is something else, and separating the two is the starting point of this unit: a program is static text, while a process is a dynamic, stateful entity.

\begin{definition}[Process]
  A \emph{process} is an abstraction of a running program. It owns
  \begin{itemize}
    \item all the memory allocated to it,
    \item all the files it has opened, and
    \item accounting information: owner, running time, memory usage, \dots
  \end{itemize}
\end{definition}

So a process is not just ``code that is running''; it is the full context the kernel uses to manage that running instance. In Linux that context is a structure called \ts{} (\Cref{sec:kernel}). The process is the most central concept in an operating system. Before studying it, we follow a C source file all the way to something that can be executed.

\subsection{Life of a C Program}

The build pipeline has four stages, each with its own tool (\Cref{fig:pipeline}). \code{gcc hello.c -o hello} runs all four at once, but each can be run separately.

\begin{figure}[htbp]
\centering
\begin{tikzpicture}[node distance=12mm]
  \node[box] (c) {\code{hello.c}};
  \node[box, right=of c] (i) {expanded code};
  \node[box, right=of i] (s) {\code{hello.s}};
  \node[box, right=of s] (o) {\code{hello.o}};
  \node[box, right=of o, fill=kspace] (e) {\code{hello}};
  \draw[arr] (c) -- node[above, tool]{cpp} (i);
  \draw[arr] (i) -- node[above, tool]{gcc} (s);
  \draw[arr] (s) -- node[above, tool]{as} (o);
  \draw[arr] (o) -- node[above, tool]{ld} (e);
  \node[box, below=9mm of $(o)!0.5!(e)$, font=\scriptsize] (lib) {\code{lib*.a} / \code{lib*.so}};
  \draw[arr] (lib) -- ($(o)!0.5!(e)$);
  \foreach \n/\t in {c/source, s/assembly, o/object file, e/executable}
    \node[below=0.5mm of \n, font=\scriptsize, color=cmt] {\t};
\end{tikzpicture}
\caption{The build pipeline: preprocessor, compiler, assembler, linker (after slide 3).}
\label{fig:pipeline}
\end{figure}

\subsubsection{The Preprocessor}

The preprocessor (\code{cpp}) expands \emph{directives} such as \code{\#include}, \code{\#define}, and \code{\#ifdef}. It performs pure text substitution and knows nothing about C semantics. \code{gcc -E main.c} runs only this stage: \code{\#include "header.h"} is replaced by the whole header, and a macro call like \code{MAX(1, 2)} is expanded in place to \code{((1) > (2) ? (1) : (2))}. Because expansion is textual, macro parameters must be parenthesized everywhere.

\subsubsection{The Compiler}

The compiler takes the expanded C code, checks the syntax, and emits assembly (\code{gcc}) or LLVM IR (\code{clang}). It also optimizes the code, sometimes dramatically.

\begin{example}[Optimization Can Erase a Function]
  Consider
\begin{lstlisting}
int main() {
  int x = 1;
  x = x - 1;
  return x;
}
\end{lstlisting}
  \noindent
  \begin{minipage}[t]{0.47\linewidth}
  \centering\small\code{gcc -S}
\begin{lstlisting}[style=sh]
pushq   %rbp
movq    %rsp, %rbp
movl    $1, -4(%rbp)
subl    $1, -4(%rbp)
movl    -4(%rbp), %eax
popq    %rbp
ret
\end{lstlisting}
  \end{minipage}\hfill
  \begin{minipage}[t]{0.47\linewidth}
  \centering\small\code{gcc -S -O2}
\begin{lstlisting}[style=sh]
xorl    %eax, %eax
ret
\end{lstlisting}
  \end{minipage}

  Without optimization, the compiler faithfully pushes, stores, subtracts, reloads, pops, and returns. With \code{-O2} it sees that the function always returns $0$ and emits two instructions.
\end{example}

\begin{keypoint}
  The code you write and the instructions the CPU executes need not correspond. Shifted line numbers and ``optimized out'' variables when debugging optimized code both come from this.
\end{keypoint}

\subsubsection{The Assembler}

The assembler (\code{as}) converts assembly into an \emph{object file}. The object file already contains machine code, but it is \emph{not yet executable}: it still has unresolved symbols. Where \code{printf} lives, for instance, is not known yet.

\subsubsection{The Linker}

The linker (\code{ld}) combines object files and libraries into the final executable.

\begin{definition}[Library]
  A \emph{library} is a collection of functions and variables. It comes in two forms:
  \begin{center}
  \begin{tabular}{llll}
    \toprule
     & a.k.a. & Created with & Extension \\
    \midrule
    static library & archive & \code{ar -rcs} & \code{.a} \\
    dynamic library & shared object & \code{gcc -shared} & \code{.so} / \code{.dylib} / \code{.dll} \\
    \bottomrule
  \end{tabular}
  \end{center}
  The three dynamic-library extensions are those of Linux, macOS, and Windows.
\end{definition}

\subsection{Static and Dynamic Linking}

The two kinds of library correspond to two ways of linking, which differ in \emph{where the library code ends up} (\Cref{fig:linking}).
\begin{itemize}
  \item \emph{Static linking.} The final program embeds every library function it uses. It is self-contained and needs no external files at run time; the cost is a larger binary, and a library upgrade requires relinking.
  \item \emph{Dynamic linking.} The final program only records ``I need library X''. The OS maps the library in when the program runs. Many programs share one copy, saving memory and disk, and a library bug fix benefits all of them at once.
\end{itemize}
\code{ldd /bin/python3} lists the dynamic libraries a program depends on.

\begin{figure}[htbp]
\centering
\begin{tikzpicture}[
    file/.style={draw, minimum width=10mm, minimum height=10mm, font=\scriptsize\ttfamily, inner sep=1pt},
    linker/.style={draw, ellipse, fill=uspace, minimum width=18mm, minimum height=9mm}]
  \node[file, fill=kspace] (o1) at (0,0.65) {hello.o};
  \node[file, fill=kspace!40] (a) at (0,-0.65) {lib*.a};
  \node[linker] (l1) at (2.3,0) {Linker};
  \node[file, fill=childc!40, minimum height=19mm, minimum width=12mm] (h1) at (4.5,0) {hello};
  \draw[arr] (o1) -- (l1); \draw[arr] (a) -- (l1); \draw[arr] (l1) -- (h1);
  \node[anchor=west, align=left] at (5.3,0)
    {\textbf{Static linking}\\The final program embeds all library functions used.};
  \draw[dashed, cmt] (-0.8,-1.45) -- (13.2,-1.45);
  \node[file, fill=kspace] (o2) at (0,-2.25) {hello.o};
  \node[file, fill=kspace!40] (so) at (0,-3.55) {lib*.so};
  \node[linker] (l2) at (2.3,-2.9) {Linker};
  \node[file, fill=childc!40] (h2) at (4.5,-2.9) {hello};
  \draw[arr] (o2) -- (l2); \draw[arr] (so) -- (l2); \draw[arr] (l2) -- (h2);
  \node[anchor=west, align=left] at (5.3,-2.9)
    {\textbf{Dynamic linking}\\The final program does not embed library functions;\\the OS loads dynamic libraries when the program runs.};
\end{tikzpicture}
\caption{Static vs.\ dynamic linking (after slide 11).}
\label{fig:linking}
\end{figure}

The rest of the unit first studies the system calls a programmer uses to manage processes (create, execute, wait, signal) and then turns to the kernel to see what each one does inside.

\subsection{(SUPPLEMENT) Further Topics}
\label{sec:program-further}

\paragraph{Why an object file cannot run yet.} Call instructions in a \code{.o} file hold placeholder addresses. A \emph{relocation table} lists the spots that need real addresses, and a \emph{symbol table} lists the symbols the file defines and references. The linker resolves symbols across all object files and then patches the addresses. \code{undefined reference to 'foo'} is a symbol-resolution failure, not a compile error: it appears in the last stage of the pipeline.

\paragraph{The price of dynamic linking.} Addresses of library functions are known only at run time, so calls go through an indirect jump (the PLT and GOT). This costs a little at run time and brings ``DLL hell'': a program fails to start when the library version on the target machine does not match. Go links statically by default and Rust leans static, trading binary size for deployment determinism.

\paragraph{\code{gcc} is a driver.} \code{gcc} does not compile anything itself. It invokes \code{cpp}, \code{cc1} (the actual compiler), \code{as}, and \code{ld} as needed; \code{gcc -v hello.c} prints the commands it dispatches, which helps locate the stage that fails.
